\chapter{量子介质中的电磁响应}\label{chap:quantum-electromagnetic-media}

真空散射的平面波外态依赖平移对称性。材料、边界和开放空间改变了传播，而吸收又会把电磁能量送入未观测的自由度。单独量子化一个有损 Maxwell 算子不能给出封闭的幺正动力学，材料储库与辐射通道必须共同进入场的量子描述。

因果 Green 张量决定源产生的响应，储库算符补足耗散所伴随的量子涨落。若只保留一个正能电子工作窗，投影自能还记录被消去通道的虚跃迁；连续电子--模式耦合则须继承原外态测度。这个具体投影把微观电流顶角接到低能电子、原子及电磁环境的模型，而不改变其归一化来源。

\section{从 QED 到低能电磁模型}

完整 QED 的状态空间允许不同粒子数通道；低能实验却可能只制备并测量一个正能电子，同时用线性响应表征材料。把电流--光子耦合投影到这一电子子空间，并用因果 Green 张量承载几何、介质和耗散信息，便能明确写出电子与环境交换激发的有效耦合。

这种约化不能靠直接删除未观测自由度完成：严格消失的单电子项可以从粒子数结构判断，而虚跃迁必须通过投影自能继续保留；材料与边界则进入因果 Green 张量。这样得到的有效电子光学仍与完整 QED 使用同一电流耦合和单位约定。

Coulomb 规范中消去非动力学标势后，纵向相互作用由总电荷密度决定。若把总电荷密度按被显式保留的电子子系统与其余带电环境拆成
\[
\hat\rho_q(\vb{r})=\hat\rho_e(\vb{r})+\hat\rho_B(\vb{r}),
\]
则瞬时纵向项可写成
\begin{equation}
{
\begin{aligned}
\hat H_C={}&\frac12\int\dd^3 r\,\dd^3 r'\,
\hat\rho_e(\vb{r})V_C(\vb{r}-\vb{r}')\hat\rho_e(\vb{r}')\\
&+\frac12\int\dd^3 r\,\dd^3 r'\,
\hat\rho_B(\vb{r})V_C(\vb{r}-\vb{r}')\hat\rho_B(\vb{r}')\\
&+\int\dd^3 r\,\dd^3 r'\,
\hat\rho_e(\vb{r})V_C(\vb{r}-\vb{r}')\hat\rho_B(\vb{r}'),
\qquad
V_C(\vb{r})=\frac{1}{4\pi\varepsilon_0|\vb{r}|}.
\end{aligned}}
\label{eq:coulomb-subsystem-decomposition-dp75}
\end{equation}
前两项分别属于电子子系统和环境自身的纵向能量；量子场论中它们还必须区分正规序二体项与需要重整化的单粒子自能。第三项才是两个不同子系统之间的纵向 Coulomb 耦合，其中总表达式原有的 $1/2$ 已被两项相等的交叉乘积抵消。因而“严格一电子扇区中电子--电子二体项为零”并不意味着单个电子失去 Coulomb 相互作用：它仍可通过最后一项与外部电荷或材料电荷涨落耦合。

利用单光子交换得到的 Coulomb 核 \eqnrefs{eq:coulomb-Vq-r13,eq:coulomb-from-qed-r13}，若只保留电子扇区，正规序后的二体部分在动量表象中具有一般结构
\begin{equation}
 \hat H_C^{(2)}
 =\frac12\sum_{ss'}\int\dd^3 p\,\dd^3 p'\,\dd^3 q\,
 \widetilde V_C(\vb{q})\,
 b_s^\dagger(\vb{p}+\vb{q})
 b_{s'}^\dagger(\vb{p}'-\vb{q})
 b_{s'}(\vb{p}')b_s(\vb{p}),
 \label{eq:coulomb-two-body-momentum-r2}
\end{equation}
其中所有 $(2\pi\hbar)^3$ 归一化因子已吸收到积分测度与 $\widetilde V_C$ 的定义中。这个算符始终含两个电子湮灭算符，因此对任意严格一电子态 $|\Psi_1\rangle$ 都有


\begin{proposition}[正规序二体算子在一粒子空间上为零]
设二体算子正规序且每项包含两个电子湮灭算符，态严格属于一电子 Fock 子空间；不含外部势或完整 QED 自能。则有
\begin{equation}
 {\hat H_C^{(2)}|\Psi_1\rangle=0.}
\label{eq:one-electron-coulomb-zero-r2}
\end{equation}
\end{proposition}
\begin{proof}
任意严格一电子态可写成

\begin{equation*}
 |\Psi_1\rangle
 =\sum_r\int\dd^3 k\,f_r(\vb{k})
 b_r^\dagger(\vb{k})|0\rangle.
\end{equation*}
\eqnrefs{eq:coulomb-two-body-momentum-r2} 中的算符核包含
\begin{equation*}
 b_{s'}(\vb{p}')b_s(\vb{p})b_r^\dagger(\vb{k})|0\rangle.
\end{equation*}
先把最右侧湮灭算符移过产生算符：
\begin{align*}
 b_s(\vb{p})b_r^\dagger(\vb{k})|0\rangle
 &=\{b_s(\vb{p}),b_r^\dagger(\vb{k})\}|0\rangle
 -b_r^\dagger(\vb{k})b_s(\vb{p})|0\rangle\\
 &=(2\pi\hbar)^3\delta_{sr}\delta^{(3)}(\vb{p}-\vb{k})|0\rangle.
\end{align*}
第二个湮灭算符随后直接作用真空，故
\begin{equation*}
 b_{s'}(\vb{p}')b_s(\vb{p})b_r^\dagger(\vb{k})|0\rangle=0.
\end{equation*}
对 $r,\vb{k}$ 的任意线性叠加仍为零，于是 $\hat H_C^{(2)}|\Psi_1\rangle=0$。证明只使用“正规序二体算符含两个电子湮灭算符”和“一电子态只有一个电子产生算符”这两点；它不适用于外部电荷势、材料诱导自能或含虚正电子/光子的完整 QED 修正。
\end{proof}

这只说明正规序电子--电子二体 Coulomb 项在严格一电子扇区为零；单个电子仍可与外部电荷、边界诱导场和材料极化相互作用，完整 QED 中也仍存在电子自能。

\subsection{投影与因果 Green 张量}

完整 QED 的 Hilbert 空间包含电子、正电子、光子以及材料环境等自由度。低能电子光学只保留其中一部分时，采用 Feshbach--Schur 母\eqnrefs{eq:feshbach_schur_exact_positive_projection} 对未保留自由度作精确预解式消元。令 $P$ 投影到显式保留的低能电子--光子通道，$Q=I-P$ 包含正电子、未保留的光子/材料连续谱及其他被消去通道，则
\begin{equation*}
 H_{\rm eff}(z)=PHP+\Sigma_P(z),
 \qquad
 \Sigma_P(z)=PHQ(z-QHQ)^{-1}QHP.
\end{equation*}
复能量参数 $z$ 的国际单位制单位为 J。$\Sigma_P(z)$ 完整保留 $P\to Q\to P$ 的虚跃迁，因此约化不是简单删除 $Q$ 空间：若 $Q$ 含连续谱，边界值 $\Sigma_P(E+\ii0^+)$ 可出现虚部，对应向被消去通道的不可逆损失；若 $Q$ 主要由远失谐离散态组成，实部则产生色散能级移动。

Hilbert 空间的精确消元还需要与相互作用量子场论的 Green 函数联系。对电磁环境而言，被消去的光子和材料自由度最终通过具有正确因果边界条件的推迟 Green 张量重新进入电子动力学；因此“投影自能”与“环境 Green 张量”不是两套无关的有效理论，而是同一预解式在离散子空间与连续电磁自由度上的两种表示。

% 宏观 QED：有损介质的连续储库量子化

Maxwell 方程把线性材料响应压缩到电磁算符 \(\mathsf L_{\rm EM}\)、推迟并矢 Green 张量 \(\mathbf G^R\) 与相应光学恒等式。存在吸收时，电磁子系统本身不是封闭哈密顿系统；量子化必须显式加入环境储库，使场算符既保持正则对易关系，又让真空与热涨落由同一个耗散核 \(\operatorname{Im}\mathbf G^R\) 决定。

对只含局域电损耗的介质，先让整个外部空间具有趋于零的被动吸收，并在所有空间积分完成后取极限。这样连续储库同时涵盖材料吸收与外部辐射通道，\eqnrefs{eq:green-optical-identity-SI-r12} 的体积分形式才完整；若外部始终取严格真空，等价写法须另外加入入射自由场算符。引入连续玻色储库算符
\begin{equation*}
 \hat{\mathbf f}(\mathbf r,\omega),
 \qquad \omega>0,
\end{equation*}
并规定
\begin{equation}
 \boxed{
 [\hat f_i(\mathbf r,\omega),
 \hat f_j^\dagger(\mathbf r',\omega')]
 =\delta_{ij}\delta^{(3)}(\mathbf r-\mathbf r')\delta(\omega-\omega').}
 \label{eq:reservoir-comm-SI-r12}
\end{equation}
介质的吸收谱进入噪声电流
\begin{equation}
 \boxed{
 \hat{\mathbf J}_{N}(\mathbf r,\omega)
 =\omega
 \sqrt{\frac{\hbar\varepsilon_0}{\pi}
 \operatorname{Im}\boldsymbol\varepsilon_r(\mathbf r,\omega)}
 \cdot\hat{\mathbf f}(\mathbf r,\omega),}
 \label{eq:noise-current-SI-r12}
\end{equation}
其中矩阵平方根取半正定分支。\eqnrefs{eq:noise-current-SI-r12} 的归一化由电场正则对易关系与 Green 张量光学恒等式共同固定，而不是独立的唯象噪声强度。


把噪声电流代入 Maxwell 推迟解，得到正频电场算符
\begin{equation}
 \boxed{
 \hat{\mathbf E}^{(+)}(\mathbf r,\omega)
 =\ii\mu_0\omega^2
 \int\dd^3 s\,
 \mathbf G^R(\mathbf r,\mathbf s;\omega)
 \cdot
 \sqrt{\frac{\hbar\varepsilon_0}{\pi}
 \operatorname{Im}\boldsymbol\varepsilon_r(\mathbf s,\omega)}
 \cdot\hat{\mathbf f}(\mathbf s,\omega).}
 \label{eq:macqed-Eplus-SI-r12}
\end{equation}
完整时域场由正、负频部分叠加：
\begin{equation*}
 \hat{\mathbf E}(\mathbf r,t)
 =\int_0^\infty\dd\omega\,
 \hat{\mathbf E}^{(+)}(\mathbf r,\omega)\ee^{-\ii\omega t}
 +\mathrm{H.c.}
\end{equation*}
因此介质的经典吸收谱不仅决定平均响应，也决定量子场涨落能够从哪些空间与频率通道进入。

有限温度 \(T_{\rm th}>0\) 下，定义 Bose--Einstein 占据数
\begin{equation*}
 n_B(\omega;T_{\rm th})
 =\frac1{\exp[\hbar\omega/(k_BT_{\rm th})]-1}.
\end{equation*}
先把热态中每个模式的平均占据写成
\begin{equation}
\boxed{
 \langle\hat f_i^\dagger(\mathbf r,\omega)
 \hat f_j(\mathbf r',\omega')\rangle_{T_{\rm th}}
 =n_B(\omega;T_{\rm th})\,
 \delta_{ij}\delta^{(3)}(\mathbf r-\mathbf r')\delta(\omega-\omega').}
\label{eq:reservoir-thermal-occupation-dp68}
\end{equation}
再结合玻色对易关系，热平衡储库的两种有序相关函数为
\begin{align}
 \langle\hat f_i^\dagger(\mathbf r,\omega)
 \hat f_j(\mathbf r',\omega')\rangle_{T_{\rm th}}
 &=n_B(\omega;T_{\rm th})\,
 \delta_{ij}\delta^{(3)}(\mathbf r-\mathbf r')\delta(\omega-\omega'),\notag\\
 \langle\hat f_i(\mathbf r,\omega)
 \hat f_j^\dagger(\mathbf r',\omega')\rangle_{T_{\rm th}}
 &=[n_B(\omega;T_{\rm th})+1]\,
 \delta_{ij}\delta^{(3)}(\mathbf r-\mathbf r')\delta(\omega-\omega').
 \label{eq:reservoir-thermal-correlators-dp43}
\end{align}

\begin{proposition}[热 Gaussian 库的二点函数]
设库模式满足 CCR，状态为自由模式的 Gibbs 态，模式频率为正；正负频与热占据数均采用正文归一化。则有以下结论（原式 \eqnrefs{eq:reservoir-thermal-correlators-dp43}）：
\begin{align*}
 \langle\hat f_i^\dagger(\mathbf r,\omega)
 \hat f_j(\mathbf r',\omega')\rangle_{T_{\rm th}}
 &=n_B(\omega;T_{\rm th})\,
 \delta_{ij}\delta^{(3)}(\mathbf r-\mathbf r')\delta(\omega-\omega'),\notag\\
 \langle\hat f_i(\mathbf r,\omega)
 \hat f_j^\dagger(\mathbf r',\omega')\rangle_{T_{\rm th}}
 &=[n_B(\omega;T_{\rm th})+1]\,
 \delta_{ij}\delta^{(3)}(\mathbf r-\mathbf r')\delta(\omega-\omega').
\end{align*}
\end{proposition}
\begin{proof}
先看一个频率为 $\omega$ 的离散玻色模式。热平衡密度算符为
\begin{equation*}
 \rho_{\rm th}
 =\frac{\ee^{-\beta\hbar\omega\hat a^\dagger\hat a}}
 {\operatorname{Tr}\ee^{-\beta\hbar\omega\hat a^\dagger\hat a}}
 =(1-q)\sum_{n=0}^{\infty}q^n|n\rangle\langle n|,
 \qquad q=\ee^{-\beta\hbar\omega}.
\end{equation*}
归一化使用几何级数 $\sum_{n\ge0}q^n=(1-q)^{-1}$。因此
\begin{align*}
 \langle\hat a^\dagger\hat a\rangle_{\rm th}
 &=(1-q)\sum_{n=0}^{\infty}nq^n\\
 &=(1-q)q\frac{\dd}{\dd q}\frac{1}{1-q}
 =\frac{q}{1-q}
 =\frac{1}{\ee^{\beta\hbar\omega}-1}
 \equiv n_B(\omega;T_{\rm th}).
\end{align*}
相反次序不能简单把两个算符交换；由
$\hat a\hat a^\dagger=\hat a^\dagger\hat a+I$，
\begin{equation*}
 \langle\hat a\hat a^\dagger\rangle_{\rm th}
 =n_B+1.
\end{equation*}
同一个结果也可直接从数态看出：
$\hat a\hat a^\dagger|n\rangle=(n+1)|n\rangle$。而
$\langle\hat a\hat a\rangle=\langle\hat a^\dagger\hat a^\dagger\rangle=0$，因为热态在数基中对角，而这两个算符把 $|n\rangle$ 移到与其正交的 $|n\mp2\rangle$。

宏观 QED 储库只是把这个单模结果推广到连续标签
$\lambda=(i,\mathbf r,\omega)$。规范化对易关系
\begin{equation*}
 [\hat f_i(\mathbf r,\omega),
 \hat f_j^\dagger(\mathbf r',\omega')]
 =\delta_{ij}\delta^{(3)}(\mathbf r-\mathbf r')
 \delta(\omega-\omega')
\end{equation*}
说明不同连续模式彼此独立，因此离散的 Kronecker $\delta$ 被 Dirac $\delta$ 取代：
\begin{align*}
 \langle\hat f_i^\dagger(\mathbf r,\omega)
 \hat f_j(\mathbf r',\omega')\rangle
 &=n_B(\omega;T_{\rm th})\,
 \delta_{ij}\delta^{(3)}(\mathbf r-\mathbf r')\delta(\omega-\omega'),\\
 \langle\hat f_i(\mathbf r,\omega)
 \hat f_j^\dagger(\mathbf r',\omega')\rangle
 &=[n_B(\omega;T_{\rm th})+1]\,
 \delta_{ij}\delta^{(3)}(\mathbf r-\mathbf r')\delta(\omega-\omega').
\end{align*}
这就是正文\eqnrefs{eq:reservoir-thermal-correlators-dp43}。$T_{\rm th}\to0$ 时 $n_B\to0$，但第二行的单位项保留；它来自玻色对易代数本身，而不是来自热占据，因此对应真空涨落。
\end{proof}

第二行比第一行多出的单位项来自玻色对易关系；它区分环境吸收量子与向系统发射量子的两种有序相关函数。零温时 \(n_B\to0\)，储库仍保留这个单位项。把它代入\eqnrefs{eq:macqed-Eplus-SI-r12}，再用 Green 张量光学恒等式，可得
\begin{equation}
\boxed{
 \langle0|\hat{\mathbf E}^{(+)}(\mathbf r,\omega)
 \hat{\mathbf E}^{(-)}(\mathbf r',\omega')|0\rangle
 =\frac{\hbar\mu_0}{\pi}\omega^2
 \operatorname{Im}\mathbf G^R(\mathbf r,\mathbf r';\omega)
 \delta(\omega-\omega').}
\label{eq:macqed-vacuum-field-correlation-dp68}
\end{equation}
因此同一个 \(\operatorname{Im}\mathbf G^R\) 一方面记录材料吸收和辐射通道，另一方面固定零点电场涨落的谱权重。



若系统无损、封闭，并且 Maxwell 广义本征问题能够化为厄米形式，则推迟 Green 张量退化为离散实频率极点之和：
\begin{equation}
 \mathbf G^R(\mathbf r,\mathbf r';\omega)
 =\sum_\lambda
 \frac{\mathbf E_\lambda(\mathbf r)\otimes\mathbf E_\lambda^*(\mathbf r')}
 {\mathcal N_\lambda[\omega_\lambda^2-(\omega+\ii0^+)^2]}.
 \label{eq:G-mode-expansion-SI-r12}
\end{equation}
这里 \(\mathcal N_\lambda\) 定义为相应极点留数的电磁能量归一化。色散介质中它一般包含
\begin{equation*}
 \frac{\partial[\omega\epsilon(\omega)]}{\partial\omega},
 \qquad
 \frac{\partial[\omega\mu(\omega)]}{\partial\omega},
\end{equation*}
不能直接照搬真空中的 \(\int(\epsilon|\mathbf E|^2+\mu|\mathbf H|^2)\)。一旦存在真实吸收，离散正常模展开不再完备，必须回到连续储库与 Green 张量表示。


有损介质比无损正常模体系多出一个不可缺少的对象：连续储库。它与经典 Maxwell Green 张量并不是两套独立理论；材料吸收先进入 \(\operatorname{Im}\boldsymbol\varepsilon_r\)，随后通过\eqnrefs{eq:green-optical-identity-SI-r12} 决定 \(\operatorname{Im}\mathbf G^R\)，最终由\eqnrefs{eq:macqed-vacuum-field-correlation-dp68} 固定量子场涨落。无损封闭系统的离散正常模只是在耗散消失后的特殊极限。



连续电子与单量子模式的有效耦合必须与完整 QED 的外态归一化一致。把傅里叶归一化电子算符换成普通 Dirac-$\delta$ 归一化后，
\[
 \hat c_{\vb{p} s}
 \equiv\frac{\hat b_s(\vb{p})}{(2\pi\hbar)^{3/2}},
 \qquad
 \{\hat c_{\vb{p} s},\hat c_{\vb{p}'s'}^\dagger\}
 =\delta_{ss'}\delta^{(3)}(\vb{p}-\vb{p}'),
\]
并把模式 $\lambda$ 的一量子矢势矩阵元定义为
\[
 \vb{A}_\lambda(\vb{r})
 \equiv
 \langle0|\hat{\vb{A}}^{(+)}(\vb{r})|1_\lambda\rangle,
 \qquad
 [\vb{A}_\lambda]=\mathrm{V\,s\,m^{-1}}.
\]
模式指标 $\lambda$ 既可标记封闭系统中的离散模式，也可标记开放有损环境中的连续储库通道；不同归一化只改变相应的模式测度，不改变电子自旋子矩阵元的结构。

QED 横向相互作用哈密顿量
\[
 \hat H_{jA}
 =-q_{\mathrm e} c\int\dd^3 r:\!\hat{\bar\psi}(\vb{r})
 \boldsymbol\gamma\cdot\hat{\vb{A}}_T(\vb{r})
 \hat\psi(\vb{r})\!:
\]
在正能单电子扇区中可写成
\[
 \hat H_{jA}
 \supset
 \hbar\sum_{\lambda,ss'}\int\dd^3 p\,\dd^3 p'\,
 \mathcal G^{(\lambda)}_{\vb{p}'s',\vb{p}s}
 \hat c_{\vb{p}'s'}^\dagger\hat c_{\vb{p}s}\hat a_\lambda
 +\mathrm{H.c.},
\]
其中决定低能有效动力学的是连续核
\begin{equation}
 {
 \begin{aligned}
 \hbar\mathcal G^{(\lambda)}_{\vb{p}'s',\vb{p}s}
 ={}&-q_{\mathrm e} c\,
 \sqrt{\frac{c}{2E_{\vb{p}'}}}
 \sqrt{\frac{c}{2E_{\vb{p}}}}\\
 &\times\int\frac{\dd^3 r}{(2\pi\hbar)^3}
 \bar u_{s'}(p')
 \boldsymbol\gamma\cdot\vb{A}_\lambda(\vb{r})
 u_s(p)
 \ee^{\ii(\vb{p}-\vb{p}')\cdot\vb{r}/\hbar}.
 \end{aligned}}
 \label{eq:qed_to_mode_matrix_element_detailed}
\end{equation}
两个平方根来自正能 Dirac 场的协变外线归一化，$(2\pi\hbar)^{-3}$ 来自两次连续 Fourier 展开。因而 $\mathcal G^{(\lambda)}_{\vb{p}'s',\vb{p}s}$ 不是一个脱离测度的普通数，而是与 $\dd^3 p\,\dd^3 p'$ 和 Dirac-$\delta$ 归一化算符成对定义的连续核；只有再做波包、有限盒或离散通道投影后，才得到普通的离散耦合频率。\eqnrefs{eq:qed_to_mode_matrix_element_detailed} 因而把完整 QED 的协变矩阵元投影为一般电子--量子系统哈密顿量中的低能连续耦合核。

低能投影与环境 Green 张量把完整 QED 的电流耦合落实为可用于材料中的连续核。量子光学中的偶极率是这种场涨落在局域原子跃迁上的应用；自由传播电子还要处理波包、有限作用时间、反冲与条件读出，这些具体电子测量问题由《自由电子量子光学》继续展开。这里已经给出其所需的电磁场、储库与归一化来源。转向强相互作用时，变化的则是内部表示和规范群：电磁相位的一维空间将由非阿贝尔颜色空间取代。





\begin{derivation}[储库真空态中的电场关联]
把正频场写成

\begin{equation*}
\hat E_i^{(+)}(\mathbf r,\omega)
=\int\dd^3 s\,K_{ia}(\mathbf r,\mathbf s;\omega)
\hat f_a(\mathbf s,\omega),
\end{equation*}
其中 \(K\) 包含正文\eqnrefs{eq:macqed-Eplus-SI-r12} 中的全部数值核。真空满足
\begin{equation*}
\langle0|\hat f_a(\mathbf s,\omega)
\hat f_b^\dagger(\mathbf s',\omega')|0\rangle
=\delta_{ab}\delta^{(3)}(\mathbf s-\mathbf s')\delta(\omega-\omega').
\end{equation*}
代入两份场算符后，储库的空间 \(\delta\) 函数消去一个积分：
\begin{align*}
&\langle0|\hat{\mathbf E}^{(+)}(\mathbf r,\omega)
\hat{\mathbf E}^{(-)}(\mathbf r',\omega')|0\rangle\\
&=\mu_0^2\omega^4\frac{\hbar\varepsilon_0}{\pi}
\delta(\omega-\omega')
\int\dd^3 s\,
\mathbf G^R(\mathbf r,\mathbf s;\omega)
\cdot\operatorname{Im}\boldsymbol\varepsilon_r(\mathbf s,\omega)
\cdot(\mathbf G^R)^\dagger(\mathbf r',\mathbf s;\omega).
\end{align*}
利用 \(\mu_0\varepsilon_0=1/c^2\) 和正文\eqnrefs{eq:green-optical-identity-SI-r12}，积分等于 \((c^2/\omega^2)\operatorname{Im}\mathbf G^R\)。于是
\begin{equation*}
\langle0|\hat{\mathbf E}^{(+)}(\mathbf r,\omega)
\hat{\mathbf E}^{(-)}(\mathbf r',\omega')|0\rangle
=\frac{\hbar\mu_0}{\pi}\omega^2
\operatorname{Im}\mathbf G^R(\mathbf r,\mathbf r';\omega)
\delta(\omega-\omega'),
\end{equation*}
即正文\eqnrefs{eq:macqed-vacuum-field-correlation-dp68}。
\end{derivation}

\begin{derivation}[电磁顶角在模式基中的矩阵元]
正能电子场与其 Dirac 伴随在固定时刻分别写成
\[
\hat\psi_+(\mathbf r)
=\sum_s\int\frac{\dd^3p}{(2\pi\hbar)^{3/2}}
\sqrt{\frac{c}{2E_{\mathbf p}}}
\,u_s(p)\hat c_{\mathbf ps}
 \ee^{\ii\mathbf p\cdot\mathbf r/\hbar},
\]
\[
\hat{\bar\psi}_+(\mathbf r)
=\sum_{s'}\int\frac{\dd^3p'}{(2\pi\hbar)^{3/2}}
\sqrt{\frac{c}{2E_{\mathbf p'}}}
\,\hat c_{\mathbf p's'}^\dagger\bar u_{s'}(p')
 \ee^{-\ii\mathbf p'\cdot\mathbf r/\hbar}.
\]
取矢势正频部分
\[
\hat{\mathbf A}^{(+)}_T(\mathbf r)
=\sum_\lambda\mathbf A_\lambda(\mathbf r)\hat a_\lambda.
\]
代入
\[
\hat H_{jA}=-q_{\mathrm e} c\int\dd^3r:
\hat{\bar\psi}\,\boldsymbol\gamma\cdot\hat{\mathbf A}_T\hat\psi:
\]
并抽取算符单项
$\hat c_{\mathbf p's'}^\dagger\hat c_{\mathbf ps}\hat a_\lambda$，其系数为
\begin{align*}
K^{(\lambda)}_{\mathbf p's',\mathbf ps}
={}&-q_{\mathrm e} c
\sqrt{\frac{c}{2E_{\mathbf p'}}}
\sqrt{\frac{c}{2E_{\mathbf p}}}\\
&\times\int\frac{\dd^3r}{(2\pi\hbar)^3}
\bar u_{s'}(p')
\boldsymbol\gamma\cdot\mathbf A_\lambda(\mathbf r)
 u_s(p)
 \ee^{\ii(\mathbf p-\mathbf p')\cdot\mathbf r/\hbar}.
\end{align*}
为了证明这确实就是物理通道矩阵元，而不只是“比较系数”，用 Dirac-$\delta$ 归一化一电子态
\[
|\mathbf p,s;1_\lambda\rangle
\equiv
\hat c_{\mathbf ps}^\dagger\hat a_\lambda^\dagger|0\rangle
\]
计算
\begin{align*}
&\langle0|\hat a_{\lambda'}
\hat c_{\mathbf p_f s_f}
\hat H_{jA}^{(+)}
\hat c_{\mathbf p_i s_i}^\dagger
\hat a_\lambda^\dagger|0\rangle\\
&\qquad=
K^{(\lambda)}_{\mathbf p_f s_f,\mathbf p_i s_i}
\,\delta_{\lambda'\lambda}.
\end{align*}
这里电子部分使用
\[
\{\hat c_{\mathbf ps},\hat c_{\mathbf p's'}^\dagger\}
=\delta_{ss'}\delta^{(3)}(\mathbf p-\mathbf p')
\]
两次收缩，光子部分使用
$\langle0|\hat a_{\lambda'}\hat a_\lambda^\dagger|0\rangle=\delta_{\lambda'\lambda}$；所有没有与外态匹配的算符项都湮灭真空。因而有效哈密顿量 写成
\[
\hat H_{jA}^{(+)}
=\hbar\sum_\lambda\int\dd^3p\dd^3p'\,
\mathcal G^{(\lambda)}_{\mathbf p's',\mathbf ps}
\hat c_{\mathbf p's'}^\dagger\hat c_{\mathbf ps}\hat a_\lambda
\]
时，必须有
\[
\hbar\mathcal G^{(\lambda)}_{\mathbf p's',\mathbf ps}
=K^{(\lambda)}_{\mathbf p's',\mathbf ps},
\]
即正文\eqnrefs{eq:qed_to_mode_matrix_element_detailed}。这样 $(2\pi\hbar)^{-3}$、两个 $\sqrt{c/(2E)}$ 与一量子模式归一化都由场展开和外态收缩唯一固定，而不是后来人为补上的归一化因子。
\end{derivation}

